\section*{Problemas}

\subsection*{Enunciados de los problemas}

% Recordar
\begin{problema}
	\label{prob:ba2e465}
	Encontrar el rango y las desviaciones media y estándar de los siguientes arreglos de datos no agrupados:
	\begin{enumerate}
		\item $12, 6, 7, 3, 15, 10, 18, 5$
		\item $9, 3, 8, 8, 9, 8, 9, 18.$
	\end{enumerate}
	Compruebe sus resultados con \texttt{Python}.
\end{problema}

% Comprender
\begin{problema}
	\label{prob:c5833ad}
	Si $d=X-P$ representan las desviaciones de una variable $X$ respecto a un pivote constante $P$, demuestre que la desviación estándar es invariante ante traslaciones —es decir, que su valor no depende de qué pivote $P$ se elija— y explique por qué esta propiedad es deseable para una medida de dispersión.
\end{problema}

% Aplicar
\begin{problema}
	\label{prob:d87f55c}
	Dos grupos de estudiantes tienen las siguientes características muestrales:
	\begin{itemize}
		\item Grupo A: $n_1 = 30$ estudiantes, $\bar{X}_1 = 75$ puntos, $s_1 = 8$ puntos
		\item Grupo B: $n_2 = 20$ estudiantes, $\bar{X}_2 = 75$ puntos, $s_2 = 12$ puntos
	\end{itemize}
	Encuentre la desviación estándar del grupo combinado de $n_1 + n_2 = 50$ estudiantes.
\end{problema}

% Analizar
\begin{problema}
	\label{prob:6112df0}
	Demuestre rigurosamente que para cualquier conjunto de datos (no agrupados o agrupados):
	\begin{align}
		s &= \sqrt{\dfrac{\sum X^{2}}{N}-\left( \dfrac{\sum X}{N} \right)^{2}} = \sqrt{\overline{X^{2}}-\overline{X}^{2}}\\
		s &= \sqrt{\dfrac{\sum fX^{2}}{N}-\left( \dfrac{\sum fX}{N} \right)^{2}} = \sqrt{\overline{X^{2}}-\overline{X}^{2}}
	\end{align}
\end{problema}

% Evaluar
\begin{problema}
	\label{prob:dd18a2e}
	El coeficiente de variación se define como $CV = (s/\bar{X}) \times 100\%$. Compare e interprete la variabilidad relativa en las siguientes dos situaciones prácticas:
	\begin{itemize}
		\item Los salarios de dos empresas: Empresa A tiene $\bar{X}_A = \$45,000$ y $s_A = \$8,000$; Empresa B tiene $\bar{X}_B = \$38,000$ y $s_B = \$6,500$.
		\item Las estaturas y pesos de un grupo antropométrico: estatura media $170$ cm con $s = 8$ cm; peso medio $70$ kg con $s = 10$ kg.
	\end{itemize}
	Evalúe en cada caso cuál de las dos cantidades comparadas exhibe mayor dispersión relativa, y explique por qué la desviación estándar sola (sin dividir entre la media) sería insuficiente para esta comparación.
\end{problema}

% Crear
\begin{problema}
	\label{prob:78de91c}
	Un analista de calidad agrupa el peso (en gramos) de 50 piezas manufacturadas en la siguiente tabla, con amplitud de clase constante $c=4$ y pivote en la marca de la clase central $P=52$:
	\begin{center}
		\begin{tabular}{ccc}
			\hline
			Marca de clase $X_i$ & Frecuencia $f_i$ & $u_i=(X_i-P)/c$ \\
			\hline
			44 & 4  & $-2$ \\
			48 & 10 & $-1$ \\
			52 & 20 & $0$ \\
			56 & 12 & $1$ \\
			60 & 4  & $2$ \\
			\hline
		\end{tabular}
	\end{center}
	Diseñe la solución completa usando el método de cambio de escala y origen ($s = c \cdot s_u$): calcule $s_u$ a partir de la columna $u_i$ y úselo para obtener la desviación estándar $s$ del peso de las piezas en gramos.
\end{problema}

\subsection*{Soluciones detalladas}

% Recordar
\begin{solproblema}[prob:ba2e465]
	\begin{enumerate}
		\item Para el arreglo $12,6,7,3,15,10,18,5$:
		\begin{itemize}
			\item \textbf{Rango:} $\max - \min = 18 - 3 = 15$
			\item \textbf{Media:} $\bar{X} = \frac{12+6+7+3+15+10+18+5}{8} = \frac{76}{8} = 9.5$
			\item \textbf{Desviación media ($DM$):}
			\begin{align}
				DM &= \frac{|12-9.5| + |6-9.5| + |7-9.5| + |3-9.5| + |15-9.5| + |10-9.5| + |18-9.5| + |5-9.5|}{8} \\
				&= \frac{2.5 + 3.5 + 2.5 + 6.5 + 5.5 + 0.5 + 8.5 + 4.5}{8} = \frac{34}{8} = 4.25
			\end{align}
			\item \textbf{Desviación estándar ($s$):}
			\begin{align}
				s &= \sqrt{\frac{(12-9.5)^2 + (6-9.5)^2 + (7-9.5)^2 + (3-9.5)^2 + (15-9.5)^2 + (10-9.5)^2 + (18-9.5)^2 + (5-9.5)^2}{8}} \\
				&= \sqrt{\frac{6.25 + 12.25 + 6.25 + 42.25 + 30.25 + 0.25 + 72.25 + 20.25}{8}} = \sqrt{\frac{190}{8}} = \sqrt{23.75} \approx 4.87
			\end{align}
		\end{itemize}

		\item Para el arreglo $9,3,8,8,9,8,9,18$:
		\begin{itemize}
			\item \textbf{Rango:} $18 - 3 = 15$
			\item \textbf{Media:} $\bar{X} = \frac{9+3+8+8+9+8+9+18}{8} = \frac{72}{8} = 9$
			\item \textbf{Desviación media ($DM$):}
			\begin{align}
				DM &= \frac{|9-9| + |3-9| + |8-9| + |8-9| + |9-9| + |8-9| + |9-9| + |18-9|}{8} \\
				&= \frac{0 + 6 + 1 + 1 + 0 + 1 + 0 + 9}{8} = \frac{18}{8} = 2.25
			\end{align}
			\item \textbf{Desviación estándar ($s$):}
			\begin{align}
				s &= \sqrt{\frac{0^2 + (-6)^2 + (-1)^2 + (-1)^2 + 0^2 + (-1)^2 + 0^2 + 9^2}{8}} = \sqrt{\frac{120}{8}} = \sqrt{15} \approx 3.87
			\end{align}
		\end{itemize}
	\end{enumerate}
\end{solproblema}

\begin{lstlisting}[language=Python]
import numpy as np
# Conjunto 1
data1 = np.array([12, 6, 7, 3, 15, 10, 18, 5])
print("Conjunto 1:")
print(f"Rango: {np.ptp(data1)}")
print(f"Media: {np.mean(data1)}")
print(f"Desviación media: {np.mean(np.abs(data1 - np.mean(data1)))}")
print(f"Desviación estándar: {np.std(data1)}")
# Conjunto 2
data2 = np.array([9, 3, 8, 8, 9, 8, 9, 18])
print("\nConjunto 2:")
print(f"Rango: {np.ptp(data2)}")
print(f"Media: {np.mean(data2)}")
print(f"Desviación media: {np.mean(np.abs(data2 - np.mean(data2)))}")
print(f"Desviación estándar: {np.std(data2)}")
\end{lstlisting}

% Comprender
\begin{solproblema}[prob:c5833ad]
	Si $d_j = X_j - P$, entonces $X_j = d_j + P$. La media de $X$ en términos del pivote es:
	\begin{align}
		\bar{X} = \frac{\sum f_j X_j}{N} = \frac{\sum f_j (d_j + P)}{N} = \frac{\sum f_j d_j}{N} + P = \bar{d} + P.
	\end{align}
	La desviación respecto a la media se transforma según:
	\begin{align}
		X_j - \bar{X} = (d_j + P) - (\bar{d} + P) = d_j - \bar{d}.
	\end{align}
	El término $P$ se cancela exactamente, así que la varianza resulta idéntica sea cual sea el pivote elegido:
	\begin{align}
		s^2 = \frac{\sum f_j (X_j - \bar{X})^2}{N} = \frac{\sum f_j (d_j - \bar{d})^2}{N}.
	\end{align}
	Esta invarianza es deseable porque la dispersión de un conjunto de datos es una propiedad intrínseca de qué tan alejados están los datos entre sí, no de dónde se coloque arbitrariamente el origen de medición: trasladar todos los datos (cambiar unidades de referencia, restar una constante) no debe alterar cuánto varían entre sí.
\end{solproblema}

% Aplicar
\begin{solproblema}[prob:d87f55c]
	Como ambos grupos tienen la misma media ($\bar{X}_1 = \bar{X}_2 = 75$), la media combinada es también $\bar{X} = 75$. En este caso sin diferencia entre medias de grupo, la varianza combinada es el promedio ponderado de las varianzas individuales:
	\begin{align}
		s^2 &= \frac{n_1 s_1^2 + n_2 s_2^2}{n_1 + n_2} = \frac{30 \cdot 8^2 + 20 \cdot 12^2}{30 + 20} \\
		&= \frac{30 \cdot 64 + 20 \cdot 144}{50} = \frac{1920 + 2880}{50} = \frac{4800}{50} = 96
	\end{align}
	Por lo tanto, la desviación estándar del grupo combinado es:
	\begin{align}
		s = \sqrt{96} \approx 9.80 \text{ puntos.}
	\end{align}
\end{solproblema}

% Analizar
\begin{solproblema}[prob:6112df0]
	Partimos de la definición formal de varianza muestral (o poblacional dividida entre $N$):
	\begin{align}
		s^2 &= \frac{\sum (X_j - \bar{X})^2}{N} = \frac{\sum (X_j^2 - 2X_j\bar{X} + \bar{X}^2)}{N} \\
		&= \frac{\sum X_j^2 - 2\bar{X}\sum X_j + N\bar{X}^2}{N} = \frac{\sum X_j^2}{N} - 2\bar{X}\frac{\sum X_j}{N} + \bar{X}^2 \\
		&= \frac{\sum X_j^2}{N} - 2\bar{X}\cdot\bar{X} + \bar{X}^2 = \frac{\sum X_j^2}{N} - \bar{X}^2 = \overline{X^2} - \overline{X}^2.
	\end{align}
	Tomando la raíz cuadrada a ambos lados, se demuestra la identidad:
	\begin{align}
		s = \sqrt{\overline{X^2} - \overline{X}^2}.
	\end{align}
	Para datos con frecuencias, el mismo razonamiento algebraico es válido al reemplazar $\sum X_j^2$ por $\sum f_j X_j^2$ y $\sum X_j$ por $\sum f_j X_j$, ya que $N=\sum f_j$ y $\sum f_j$ actúa como el número total de observaciones ponderadas.
\end{solproblema}

% Evaluar
\begin{solproblema}[prob:dd18a2e]
	\textbf{Salarios:}
	\begin{align}
		CV_A &= \frac{8000}{45000} \times 100\% \approx 17.78\% \\
		CV_B &= \frac{6500}{38000} \times 100\% \approx 17.11\%
	\end{align}
	Ambas empresas muestran una dispersión salarial relativa prácticamente similar (alrededor del 17\%), aunque la Empresa A tiene una variabilidad ligeramente mayor en relación con el monto de su nómina promedio. Comparar directamente $s_A=\$8{,}000$ contra $s_B=\$6{,}500$ sería engañoso, porque ignora que la Empresa A también paga salarios promedio más altos.

	\textbf{Estaturas vs. pesos:}
	\begin{align}
		CV_{\text{estatura}} &= \frac{8}{170} \times 100\% \approx 4.71\% \\
		CV_{\text{peso}} &= \frac{10}{70} \times 100\% \approx 14.29\%
	\end{align}
	A pesar de que las unidades físicas son distintas (cm vs.\ kg), el coeficiente de variación permite su comparación adimensional directa. Se concluye que el peso presenta el triple de variabilidad relativa que la estatura en dicha población. En ambos casos, la desviación estándar sola es insuficiente para comparar porque mezcla la escala absoluta de la variable con su dispersión: dos variables con la misma $s$ pero medias muy distintas (o medidas en unidades distintas, como cm y kg) no son comparables en términos de variabilidad relativa sin normalizar por la media.
\end{solproblema}

% Crear
\begin{solproblema}[prob:78de91c]
	Con $N=\sum f_i = 4+10+20+12+4=50$, calculamos $\sum f_iu_i$ y $\sum f_iu_i^2$:
	\begin{align}
		\sum f_iu_i &= 4(-2)+10(-1)+20(0)+12(1)+4(2) = -8-10+0+12+8 = 2 \\
		\sum f_iu_i^2 &= 4(4)+10(1)+20(0)+12(1)+4(4) = 16+10+0+12+16 = 54
	\end{align}
	La desviación estándar codificada es:
	\begin{align}
		s_u = \sqrt{\dfrac{\sum f_iu_i^2}{N} - \left(\dfrac{\sum f_iu_i}{N}\right)^2} = \sqrt{\dfrac{54}{50} - \left(\dfrac{2}{50}\right)^2} = \sqrt{1.08 - 0.0016} = \sqrt{1.0784} \approx 1.0385.
	\end{align}
	Aplicando el cambio de escala $s = c \cdot s_u$ con $c=4$:
	\begin{align}
		s = 4 \times 1.0385 \approx 4.15 \text{ gramos.}
	\end{align}
	El método codificado evita manejar directamente los valores grandes de la marca de clase (44 a 60), operando en su lugar con enteros pequeños ($-2$ a $2$).
\end{solproblema}
